\item \model{Magistral}

\paragraph*{مقدمه و پارادایم داده‌ای پس‌آموزش (\en{Post-Training Data Paradigm})}
فاز پس‌آموزش (\en{Post-Training}) در مدل‌های استدلالی خانواده \model{Magistral} با هدف ارتقای قابلیت‌های زنجیره تفکر عمیق (\en{Chain-of-Thought - CoT}) بدون اتکا به تقطیر از مدل‌های انحصاری پیشین توسعه یافته است \cite{jaech2024o1, deepseek2025r1}. در این چارچوب، دو مدل با راهبردهای داده‌ای متمایز آموزش دیده‌اند:
\begin{itemize}
    \item \textbf{مدل \model{Magistral Medium}:} مبتنی بر وزن‌های پایه \en{Mistral Medium 3} \cite{mistral2025medium3} که فرایند پس‌آموزش آن \textbf{صرفاً با یادگیری تقویتی آنلاین (\en{Pure RLVR})} و بدون هیچ‌گونه داده راه‌اندازی اولیه استدلالی (\en{Cold-Start}) به انجام رسیده است.
    \item \textbf{مدل \model{Magistral Small (24B)}:} مبتنی بر \en{Mistral Small 3} که فرایند پس‌آموزش آن شامل یک فاز تنظیم دقیق تحت نظارت (\en{SFT}) با ردپاهای استدلالی پالایش‌شده و سپس یادگیری تقویتی تکمیلی (\en{RL}) است.
\end{itemize}

در این چارچوب، یک دوگانگی بنیادین در ویژگی‌های آماری داده‌ها شناسایی شده است: الگوریتم‌های بهینه‌سازی تقویتی مبتنی بر پاداش‌های قابل راستی‌آزمایی (\en{RLVR}) نیازمند مجموعه‌داده‌ای با حجم کم، خلوص فوق‌العاده بالا و سختی بسیار زیاد هستند؛ در حالی که فاز تنظیم نظارتی (\en{SFT}) به تنوع گسترده پرامپت‌ها (\en{Prompt Diversity}) برای حفظ آنتروپی و توزیع طول استدلال نیاز دارد.

\paragraph*{خط‌لوله پالایش و گزینش داده‌های ریاضی (\en{Math Data Curation})}
داده‌های حوزه ریاضیات طی یک خط‌لوله پالایش دوسطحی شامل فیلترینگ صوری و فیلترینگ سختی مبتنی بر مدل غربال‌گری شدند:
\begin{itemize}
    \item \textbf{پالایش صوری و اعتبارسنجی قطعی (\en{Format Filtering}):} از میان مجموعه‌داده اولیه شامل ۶۹۹ هزار نمونه، تمامی مسائل اثباتی کیفی (\en{Proof-based}) و چندبخشی که ارزیابی ماشینی آن‌ها با ابزارهای قطعی میسر نبود حذف شدند. همچنین مسائل چندگزینه‌ای به مسائل گزاره‌ای با پاسخ صریح عددی یا عبارتی تبدیل شدند تا احتمال پاسخ‌دهی تصادفی به صفر برسد و امکان استفاده از موتور محاسبات نمادین \en{SymPy} برای انطباق معنایی پاسخ‌ها فراهم گردد. این مرحله حجم داده‌ها را به ۵۰۱ هزار نمونه تقلیل داد.
    \item \textbf{غربال‌گری سختی دو مرحله‌ای در منطقه گلدیلاکس (\en{Goldilocks Filtering}):} برای جلوگیری از نمونه‌های بسیار ساده (با مزیت صفر) یا غیرقابل‌حل (بدون سیگنال یادگیری)، فرایند ارزیابی دو مرحله‌ای پیاده‌سازی شد:
    \begin{enumerate}
        \item در مرحله نخست، با استفاده از مدل \en{Mistral Large 2} \cite{mistral2024large2} به ازای هر مسئله ۱۶ راه‌حل نمونه‌گیری شد و مسائل دارای نرخ موفقیت ۱۰۰٪ یا ۰٪ حذف شدند. داده‌های حاصل برای آموزش یک مدل ۲۴ میلیارد پارامتری ارزیاب سختی به کار رفت.
        \item در مرحله دوم، مدل ارزیاب ۲۴B آموزش‌دیده با \en{RL} کل مجموعه ۵۰۱ هزارتایی را مجدداً با ۱۶ پاسخ به ازای هر مسئله باز-ارزیابی کرد تا از حذف اشتباه مسائل دشوار توسط مدل عمومی جلوگیری شود.
    \end{enumerate}
    \item \textbf{پالایش بر پایه تضاد اجماع (\en{Consensus-Disagreement Filtering}):} مسائلی که در آن‌ها اکثریت پاسخ‌های تولیدی مدل بر سر یک جواب به اجماع می‌رسیدند اما با پاسخ برچسب‌گذاری‌شده (\en{Ground Truth}) تضاد داشت، به عنوان خطای برچسب مبنا شناسایی و از مجموعه‌داده حذف شدند:
    \begin{equation}
        \text{mode}(\{y_i\}_{i=1}^{16}) \neq y_{\text{GT}} \quad \wedge \quad \frac{1}{16}\sum_{i=1}^{16} \mathbb{I}(y_i = \text{mode}) \ge \tau
    \end{equation}
\end{itemize}
در پایان این فرآیند، تنها \textbf{۳۸ هزار مسئله فوق‌العاده خالص و چالش‌برانگیز} برای فاز یادگیری تقویتی گزینش گردید (جدول \ref{tab:magistral_math_filter}).

\begin{table}[htbp]
\centering
\caption{مراحل پالایش و حجم دادگان آموزش ریاضی مدل \model{Magistral}}
\label{tab:magistral_math_filter}
\resizebox{0.75\textwidth}{!}{%
\begin{tabular}{lcc}
\toprule
\textbf{مرحله فیلترینگ} & \textbf{تعداد داده‌ها} & \textbf{درصد بقا نسبت به کل} \\
\midrule
داده‌های اولیه خام (\en{Initial Raw Data}) & 699,000 & 100.0\% \\
پس از پالایش صوری و اعتبارسنجی (\en{Format Filtering}) & 501,000 & 71.67\% \\
پس از ارزیابی سختی و اجماع (\en{Difficulty Filtering}) & \textbf{38,000} & \textbf{5.43\%} \\
\bottomrule
\end{tabular}%
}
\end{table}

\paragraph*{مهندسی داده‌های کدنویسی و آزمون‌های ارزیابی (\en{Code Data Engineering})}
پیکره کدنویسی بر پایه مسائل رقابتی برنامه‌نویسی با الزامات اعتبارسنجی زیر طراحی شد:
\begin{itemize}
    \item \textbf{پالایش و اصلاح آزمون‌ها با اجماع اجرا (\en{Test-Case Agreement}):} مسائل فاقد آزمون جامع یا راه‌حل‌های مرجع حذف شدند. راه‌حل‌های معتبر روی تمامی آزمون‌ها اجرا شده و آزمون‌های مبهم حذف گردیدند. چنانچه تمامی راه‌حل‌ها روی یک آزمون با خروجی یکسان مردود می‌شدند، خروجی آزمون بر اساس اجماع راه‌حل‌ها اصلاح گردید.
    \item \textbf{تولید آزمون‌های سنتتیک (\en{Synthetic Test Generation}):} برای مسائلی که پوشش تست محدودی داشتند، موارد آزمون تکمیلی به‌صورت خودکار تولید و صحت‌سنجی شدند.
    \item \textbf{تکثیر زبانی مسائل (\en{Language Duplication}):} مسائل به‌گونه‌ای تنظیم شدند که پیاده‌سازی را به دو زبان پرکاربرد \en{Python} و \en{C++20} (با استفاده از هدر از پیش کامپایل‌شده \en{bits/stdc++.h}) بطلبند. این فرایند مجموعه‌ای متشکل از \textbf{۳۵ هزار مسئله برنامه‌نویسی} تمیز را شکل داد.
\end{itemize}

\paragraph*{چندزبانگی و راهبرد مهار جابه‌جایی زبانی (\en{Multilingual Data Strategy})}
به منظور جلوگیری از پدیده اختلاط ناخواسته زبان‌ها (\en{Code-Switching}) در طول زنجیره تفکر:
\begin{itemize}
    \item \textbf{ترجمه متوازن مسائل:} ۱۰ درصد از کل مسائل ریاضی و کدنویسی انگلیسی به ۶ زبان فرانسوی، اسپانیایی، ایتالیایی، آلمانی، چینی و روسی ترجمه شد.
    \item \textbf{پاداش‌دهی بر پایه طبقه‌بند \en{FastText}:} با بهره‌گیری از طبقه‌بند زبانی \cite{joulin2016fasttext}، سه‌گانه صورت مسئله ($Q$)، افکار درونی ($C$) و پاسخ نهایی ($A$) پس از حذف نشانه‌های ساختاری و کدهای برنامه‌نویسی اعتبارسنجی شد:
    \begin{equation}
        R_{\text{lang}}(Q, C, A) = \begin{cases} 
            +0.1, & \text{if } \mathcal{C}_{\text{FastText}}(Q) = \mathcal{C}_{\text{FastText}}(C) = \mathcal{C}_{\text{FastText}}(A) \\ 
            0, & \text{otherwise} 
        \end{cases}
    \end{equation}
\end{itemize}

\paragraph*{داده‌های راه‌اندازی سرد (\en{Cold-Start SFT Data}) برای مدل کوچک}
برای مدل \model{Magistral Small (24B)}، فاز راه‌اندازی اولیه با ترکیب داده‌های زیر اجرا گردید:
\begin{itemize}
    \item \textbf{ردپاهای استخراج‌شده از \model{Magistral Medium}:} راه‌حل‌های صحیح ثبت‌شده در طول آموزش \en{RL} مدل بزرگ‌تر، با فیلتر کردن گام‌های ابتدایی و کوتاه، گزینش شدند. برای جلوگیری از سوگیری به سمت مسائل ساده، تعداد تولیدها به ازای هر مسئله محدود شد و مسائل دشوارتر بیش‌نمونه‌گیری (\en{Upsample}) شدند.
    \item \textbf{پرامپت‌های متنوع خارجی:} استفاده از مخازن متنی متنوع شامل مجموعه‌های \en{OpenThoughts} \cite{guha2025openthoughts} و \en{OpenR1} \cite{openr12025, penedo2025codeforces} که راه‌حل‌های آن‌ها توسط \model{Magistral Medium} تولید و غربال شدند.
    \item \textbf{داده‌های عمومی دستورالعمل (\en{General Alignment}):} تزریق ۱۰ درصد داده‌های مکالمه عمومی و پیروی از دستورات به پیکره \en{SFT} جهت حفظ توانمندی‌های غیر استدلالی نظیر فراخوانی ابزار (\en{Tool Calling}).
\end{itemize}
مدل پایه \en{Mistral Small 3 Instruct} برای مدت ۴ دوره (\en{Epochs}) روی این ترکیب آموزش دید و سپس وارد فاز \en{RL} شد.

\paragraph*{ردپاهای متن‌باز در مقیاس بزرگ (\en{Large-Scale OSS Traces})}
در یک سناریوی پژوهشی مستقل، مدل \en{Mistral Medium 3} با ۱.۳ میلیون ردپای استدلالی برآمده از مدل \en{DeepSeek-R1} (برگرفته از \en{OpenThoughts} و \en{OpenR1}) تنظیم نظارتی گردید. سپس یادگیری تقویتی صرفاً روی دشوارترین زیرمجموعه مسائل اعمال شد که موجب جهش بیش از ۱۰ واحد درصدی در بنچمارک \en{AIME'25} و ۵ واحد درصدی در \en{LiveCodeBench} شد و سطح عملکردی هم‌تراز با مدل‌های پیشرو ایجاد نمود.

\paragraph*{تدریج داده‌ای و پویایی طول بافت (\en{Curriculum \& Length Dynamics})}
فرآیند آموزش تقویتی به شیوه برنامه درسی مرحله‌ای (\en{Curriculum}) مدیریت شد:
\begin{itemize}
    \item \textbf{پویایی سختی داده‌ها:} همگام با تسلط مدل، مسائل کاملاً حل‌شده از چرخه خارج و مسائل دشوارتر جایگزین گردیدند.
    \item \textbf{برنامه درسی طول تولید:} سقف طول تولید بدون جریمه ($l_{\max} - l_{\text{cache}}$) در سه گام متوالی از $16\text{K}$ به $24\text{K}$ و سپس به $32\text{K}$ توکن افزایش یافت.
    \item \textbf{مقیاس دسته‌ها:} به منظور مدیریت بار حافظه جدول \en{KV-Cache} همگام با افزایش طول توالی‌ها، اندازه دسته دنباله‌ها ($n_{\text{batch}}$) در سه مرحله از $8\text{K}$ به $4\text{K}$ و سپس به $2\text{K}$ کاهش داده شد.
\end{itemize}

\paragraph*{فرمولاسیون ریاضی راستی‌آزمایی و پاداش‌دهی داده‌ها (\en{Verifiable Reward Formulation})}
در خط لوله بهینه‌سازی سیاست گروهی نسبی (\en{GRPO}) \cite{shao2024deepseekmath, yu2025dapo}، امتیاز هر نمونه بر اساس رابطه زیر تخصیص می‌یابد:
\begin{equation}
    r_i = R_{\text{format}} + R_{\text{correctness}} + R_{\text{lang}} + R_{\text{length}}
\end{equation}
که در آن پاداش قالب‌بندی $R_{\text{format}} \in \{0, 0.1\}$ شرط ورود به تصحیح پاسخ است. پاداش درستی $R_{\text{correctness}} \in \{0, 0.9\}$ در صورت موفقیت کامل محاسباتی یا عبور از آزمون‌های کدنویسی اعطا می‌شود و تابع پیوسته جریمه طول به فرم زیر محاسبه می‌گردد:
\begin{equation}
    R_{\text{length}}(y) = \begin{cases} 
        0, & |y| \le l_{\max} - l_{\text{cache}} \\ 
        -0.1 \cdot \frac{|y| - (l_{\max} - l_{\text{cache}})}{l_{\text{cache}}}, & l_{\max} - l_{\text{cache}} < |y| \le l_{\max} \\ 
        -0.1, & |y| > l_{\max} 
    \end{cases}
\end{equation}

\paragraph*{تعمیم‌پذیری بین‌حوزه‌ای و چندوجهی ناشی از داده‌های متنی (\en{Data Generalization})}
تحلیل رفتار مدل آموزش‌دیده با داده‌های تفکیک‌شده حقایق تجربی مهمی را آشکار ساخت:
\begin{itemize}
    \item \textbf{انتقال بین ریاضی و کد (\en{Cross-Domain Transfer}):} مدل ۲۴B آموزش‌دیده صرفاً با داده‌های ریاضی توانست عملکرد خود در بنچمارک کدنویسی \en{LiveCodeBench} را از ۲۲.۷٪ به ۳۸.۳٪ (۱۵.۶+ واحد درصد ارتقای برون‌دامنه) برساند؛ همچنین آموزش صرفاً با داده‌های کد، دقت ریاضی در \en{AIME'24} را ۱۷.۵+ واحد درصد بهبود بخشید (جدول \ref{tab:magistral_cross_domain}).
    \item \textbf{ارتقای رایگان درک چندوجهی (\en{Multimodal Free Lunch}):} علی‌رغم این‌که فرایند \en{RL} تنها بر روی داده‌های متنی محض اجرا شد، مدل به واسطه گسترش زنجیره تفکر توانست عملکرد در آزمون‌های استدلال چندوجهی نظیر \en{MMMU} را از ۶۵.۰٪ به ۷۰.۰٪ و در \en{MMMU-Pro (Vision)} را از ۳۹.۷٪ به ۵۲.۱٪ (۱۲.۴+ واحد درصد ارتقا) برساند \cite{yue2024mmmu, yue2025mmmupro}.
\end{itemize}

\begin{table}[htbp]
\centering
\caption{تعمیم‌پذیری بین‌حوزه‌ای در آموزش‌های تک‌دامنه‌ای مدل \model{Mistral Small 24B}}
\label{tab:magistral_cross_domain}
\resizebox{0.75\textwidth}{!}{%
\begin{tabular}{lcc}
\toprule
\textbf{پیکربندی آموزش داده} & \textbf{\en{AIME'24 (Pass@1)}} & \textbf{\en{LiveCodeBench v5 (Pass@1)}} \\
\midrule
چک‌پوینت اولیه (\en{Starting Checkpoint}) & 32.2\% & 22.7\% \\
آموزش صرفاً با داده‌های ریاضی (\en{RL Math only}) & \textbf{62.5\%} & \textbf{38.3\%} (+15.6\%) \\
آموزش صرفاً با داده‌های کدنویسی (\en{RL Code only}) & \textbf{49.7\%} (+17.5\%) & \textbf{42.7\%} \\
\bottomrule
\end{tabular}%
}
\end{table}

\paragraph*{رویکردهای ناموفق در مهندسی پاداش داده‌ها (\en{Negative Results})}
دو راهبرد پیشنهادی در فاز پردازش پاداش داده‌ها با شکست مواجه شدند:
\begin{itemize}
    \item \textbf{پاداش متناسب برای آزمون‌های کدنویسی (\en{Proportional Rewards}):} اعطای پاداش بر اساس نسبت تست‌های پاس‌شده منجر به افت ۲ درصدی در بنچمارک \en{LiveCodeBench} و مهار رشد طول زنجیره تفکر شد؛ چرا که سیگنال‌های کاذب به کدهای واجد خطاهای الگوریتمی اختصاص می‌یافت.
    \item \textbf{جریمه آنتروپی در تلفات (\en{Entropy Bonus}):} افزودن ترم آنتروپی به تابع زیان پایداری آموزش را برهم زد (افت شدید آنتروپی در ریاضی و انفجار آنتروپی در کد). در مقابل، تنظیم مستقیم حد بالای برش $\varepsilon_{\text{high}}$ در بازه $[0.26, 0.28]$ پایداری بهینه‌سازی را تضمین کرد.
\end{itemize}

\paragraph*{ارزیابی آماری جامع پیکربندی‌های داده‌ای پس‌آموزش}
جدول \ref{tab:magistral_post_train_eval} عملکرد نهایی مدل‌ها تحت استراتژی‌های گوناگون داده‌های پس‌آموزش را مقایسه می‌کند. این نتایج اثبات می‌کنند که مدل \model{Magistral Medium} صرفاً با اتکا بر ۳۸ هزار داده ریاضی و ۳۵ هزار داده کدنویسی به جهش چشمگیر در \en{AIME'24} (از ۲۶.۸٪ به ۷۳.۶٪ در سناریوی \en{Pass@1} و ۹۰.۰٪ در سناریوی \en{Maj@64}) دست یافته است.

\begin{table}[htbp]
\centering
\caption{مقایسه عملکردی مدل‌های \model{Magistral} بر اساس پیکربندی‌های مختلف داده‌های پس‌آموزش}
\label{tab:magistral_post_train_eval}
\resizebox{\textwidth}{!}{%
\begin{tabular}{lcccccc}
\toprule
\textbf{مدل} & \textbf{ترکیب داده \en{SFT}} & \textbf{ترکیب داده \en{RL}} & \textbf{\en{AIME'24 (Pass@1)}} & \textbf{\en{AIME'25 (Pass@1)}} & \textbf{\en{MATH-500}} & \textbf{\en{LCB v5}} \\
\midrule
\en{Mistral Medium 3} & داده عمومی چت & فاقد \en{RL} & 26.8\% & 21.2\% & 91.0\% & 29.1\% \\
\model{Magistral Medium} & فاقد \en{SFT} & 38k ریاضی + 35k کد & \textbf{73.6\%} & \textbf{64.9\%} & \textbf{94.3\%} & \textbf{59.4\%} \\
\midrule
\en{Mistral 24B (SFT only)} & ردپاهای Medium + مخازن وب & فاقد \en{RL} & 65.4\% & 55.6\% & 93.2\% & 52.2\% \\
\en{Mistral 24B (RL only)} & فاقد \en{SFT} & 38k ریاضی + 35k کد & 65.8\% & 51.9\% & 95.4\% & 46.4\% \\
\model{Magistral Small (24B)} & ردپاهای Medium + عمومی (4 دوره) & 38k ریاضی + 35k کد & \textbf{70.7\%} & \textbf{62.8\%} & \textbf{95.9\%} & \textbf{55.8\%} \\
\bottomrule
\end{tabular}%
}
\end{table}